
\subsubsection{2.1 Multi-Dimensional Trustworthiness Evaluation Frameworks}

Modern LLM evaluation has progressed from single-task accuracy benchmarks toward comprehensive multi-dimensional assessment. TrustLLM \citep{Sun2024TrustLLM} provides the most complete current framework, evaluating 16 prominent LLMs across 6 trustworthiness dimensions (truthfulness, safety, fairness, adversarial robustness, privacy, machine ethics) using a suite of sub-benchmarks. HELM \citep{Liang2022} offers a complementary perspective with 7 metrics across 30 models. MultiTrust \citep{Zhang2024MultiTrust} extends the analysis to multimodal LLMs across 5 dimensions. Crucially, all three frameworks present dimension scores and visual radar charts but do not compute the pairwise correlation structure between dimensions or test whether dimensions are statistically independent across models.

Liu et al. [2023] survey trustworthiness across 7 categories and explicitly identify cross-dimension correlation analysis as future work. Our paper delivers that analysis using the TrustLLM dataset.

\subsubsection{2.2 Trustworthiness Tradeoffs: Specific Pairs}

Several studies have identified specific trustworthiness tradeoffs without characterizing the full correlation structure. AQUA-LLM [Güngör et al., 2025] demonstrates a quantitative accuracy-robustness tradeoff under quantization and adversarial perturbation --- corroborating the directional safety-robustness tension we observe. Know Thy Judge \citep{Eiras2025KnowThyJudge} shows that RLHF-trained safety judges are brittle to style shifts, with false negative rates jumping by up to 0.24 --- consistent with our finding that RLHF-optimized models may not transfer safety behaviors to adversarial inputs. Wang et al. [2025] survey reasoning LLMs and find that chain-of-thought models show worse safety and privacy performance than standard models at equivalent scale, suggesting inverse scaling relationships for specific dimensions.

These studies confirm that trustworthiness dimensions are not uniformly correlated --- some pairs trade off while others co-move. However, they focus on specific pairs rather than the full 15-pair correlation matrix, and none controls for the scale and RLHF confounds that conflate model quality with trustworthiness geometry.

\subsubsection{2.3 Benchmark Correlation Analysis and Evaluation Compression}

The question of whether benchmark scores are redundant is well-studied for general capability benchmarks. The Epoch AI benchmark correlation study finds a median Spearman $\rho$ = 0.73 across 17 capability benchmarks, suggesting substantial redundancy in capability measurement. A PCA study [arXiv 2603.00394] decomposes benchmark scores into two principal components explaining 97.4\% of variance, with TruthfulQA loading orthogonally to the first component --- providing early evidence that truthfulness is a distinct dimension.

Minimum spanning tree methods for correlation-based redundancy analysis have been applied to financial return correlation matrices \citep{Tumminello2007MST} and ecological networks, where bootstrap stability metrics characterize the reliability of inferred tree topology. We apply this methodology to LLM trustworthiness evaluation for the first time.

\textbf{Methodological note on the Epoch AI baseline.} We reference the Epoch AI capability baseline (median $\rho$ = 0.73) as a field reference point, not a direct comparison. The Epoch AI study reports \textit{raw} Spearman correlations among capability benchmarks, while the correlations we report are \textit{partial} Spearman after confound removal. These quantities are methodologically distinct: partial correlations differ structurally from raw correlations when confounds are present, as illustrated by the safety--privacy pair (raw $\rho$ = 0.73 vs. partial $\rho$ = 0.971). Accordingly, the Epoch AI baseline establishes the scale of benchmark redundancy in the broader literature, not an equivalent comparison with our results.

\subsubsection{2.4 RLHF Effects on Trustworthiness}

The RLHF alignment paradigm is the primary mechanism through which modern LLMs are adapted for safe and helpful behavior \citep{Ouyang2022InstructGPT}. Li et al. [2025] (ICLR Oral) find that increasing RLHF does not automatically guarantee trustworthiness across all dimensions --- safety can improve while other dimensions stagnate or regress. Our findings are consistent with Li et al.: we confirm that RLHF jointly optimizes safety and ethics (via LLaMA-2 within-family evidence) while robustness follows an independent path.

\subsubsection{2.5 Positioning}

\begin{table}[htbp]
\centering
\resizebox{\columnwidth}{!}{%
\begin{tabular}{lll}
\toprule
Prior Work & Data Provided & What Is Missing \\
\midrule
TrustLLM [Sun et al., 2024] & 16-model $\times$ 6-dimension scores & Pairwise correlation analysis; no partial correlation \\
HELM [Liang et al., 2022] & 30-model $\times$ 7-metric scores & Cross-metric correlation matrix; no statistical clustering \\
Liu et al. [2023] & 7-category survey & Explicitly lists cross-dim correlation as future work \\
AQUA-LLM [Güngör et al., 2025] & Accuracy-robustness tradeoff & Single pair; no full matrix; no confound control \\
Epoch AI study & Capability benchmark correlations (raw $\rho$ = 0.73) & Capability, not trustworthiness; raw, not partial correlation \\
\bottomrule
\end{tabular}
}
\end{table}

Our work fills the gap at the intersection: the first pairwise partial Spearman correlation matrix for LLM trustworthiness, the first cluster analysis of the correlation geometry, and the first MST-based minimum evaluation set --- all using publicly available TrustLLM data.

